\documentclass[11pt]{article}
\usepackage{fontspec}
\usepackage[french,english]{babel}
\usepackage{amsmath}
\usepackage{amssymb}
\usepackage{amsthm}
\usepackage{mathrsfs}
\usepackage{tikz-cd}
\usepackage{enumitem}

\newcommand{\printedpage}[1]{\marginpar{\footnotesize p.~#1}}
\newcommand{\msp}[1]{\marginpar{\footnotesize [#1]}}
\newcommand{\Asph}{\mathrm{Asph}}
\newcommand{\Cof}{\mathrm{Cof}}

\begin{document}

\printedpage{175}
% head: A. Grothendieck : Dérivateurs, Ch. XII, Édition des Universités de Montpellier II et Paris VII
% foot: Version 30 mars 2003

% dangling arrows (from X_2, X_3, Y_2, Y_3, Z_2, Z_3) as printed
\[
\begin{tikzcd}[ampersand replacement=\&, column sep=small, row sep=small]
 \& \& {} \& \& {} \& \& U \arrow[dd, "\mu"] \\
 \& X_2 \arrow[rr] \arrow[dd] \arrow[ur, shorten >=1ex] \& \& X_3 \arrow[rr] \arrow[dd] \arrow[ur, shorten >=1ex] \& \& T \arrow[ur] \arrow[dd] \& \\
X_0 \arrow[rr, crossing over] \arrow[ur] \arrow[dd] \& \& X_1 \arrow[rr, crossing over] \arrow[ur] \arrow[dd, crossing over] \& \& S \arrow[ur, "\lambda"] \arrow[dd, crossing over] \& \& U' \arrow[dd] \\
 \& Y_2 \arrow[rr] \arrow[dd] \arrow[ur, shorten >=1.5ex] \& \& Y_3 \arrow[rr] \arrow[dd] \arrow[ur, shorten >=1.5ex] \& \& T' \arrow[ur] \arrow[dd] \& \\
Y_0 \arrow[rr, crossing over] \arrow[ur] \arrow[dd] \& \& Y_1 \arrow[rr, crossing over] \arrow[ur] \arrow[dd, crossing over] \& \& S' \arrow[ur, "\lambda'"] \arrow[dd, crossing over] \& \& U'' \\
 \& Z_2 \arrow[rr] \arrow[ur, shorten >=1.5ex] \& \& Z_3 \arrow[rr] \arrow[ur, shorten >=1.5ex] \& \& T'' \arrow[ur] \& \\
Z_0 \arrow[rr] \arrow[ur] \& \& Z_1 \arrow[rr] \arrow[ur] \& \& S'' \arrow[ur, "\lambda''"] \& \&
\end{tikzcd}
\]

The condition that $X$ be $W$-cocartesian is expressed by $\lambda \in W$, or again by $U \in \Asph_W$ (for $U$ identifies with the homotopy cofiber $Z(\lambda)$ of $\lambda$ \dots\,), likewise $Y$ $W$-cocartesian is expressed by $U' \in \Asph_W$, and $Z$ $W$-cocartesian by $U'' \in \Asph_W$. Now $U'' \sim Z(\mu)$, where $\mu : U \longrightarrow U'$ is in $W$ if $U$, $U'$ are $W$-aspheric, hence $U'' = Z(\mu)$ is then $W$-aspheric. This proves the lemma.

This shows that, assuming W(7 bis) (in fact, we have used only the ``universal'' part of W(7 bis)), namely that in
\begin{tikzcd}[ampersand replacement=\&, column sep=small, row sep=small, baseline=(\tikzcdmatrixname-2-1.base)]
T_0 \arrow[rr] \arrow[dr, "i"'] \& \& T_2 \arrow[dr, "i'"] \& \\
 \& T_1 \arrow[rr] \& \& T_3
\end{tikzcd}
a cocartesian diagram of open immersions, $i \in W \Longrightarrow i' \in W$ (instead of $\Longleftrightarrow$),

\msp{page 270}

the validity of W(7) reduces to the case where the square $Y$ is the trivial square
\begin{tikzcd}[ampersand replacement=\&, column sep=small, row sep=small, baseline=(\tikzcdmatrixname-2-1.base)]
 \& e \arrow[dr] \& \\
e \arrow[ur] \arrow[dr] \& \& e, \\
 \& e \arrow[ur] \&
\end{tikzcd}
$e$ the final object of Cat.

We can now conclude: We have (assuming the forward part of W(7 bis))
\[
\mathrm{a}' \ \& \ \mathrm{b}' \Longrightarrow \mathrm{c}',
\]
for if $Z_0$, $Z_1$ are $\in \Asph_W$, then $i_1 \in W$, hence $j_1 \in W$, and if in addition $Z_2 \in \Asph_W$, we therefore have $Z_3 \in \Asph_W$.
\[
\begin{tikzcd}[ampersand replacement=\&, column sep=small, row sep=small]
 \& Z_1 \arrow[dr, "j_2"] \& \\
Z_0 \arrow[ur, "i_1"] \arrow[dr, "i_2"'] \& \& Z_3 \\
 \& Z_2 \arrow[ur, "j_1"'] \&
\end{tikzcd}
\]

If we also have the reverse (``unusual'') part of W(7 bis), then
\[
\mathrm{b}' \ \& \ \mathrm{c}' \Longrightarrow \mathrm{a}',
\]
for if $Z_1$ and $Z_3$ are $W$-aspheric, then $j_2 \in W$, hence $i_2 \in W$, and if in addition $Z_2$ is $W$-aspheric, we conclude $Z_0$ $W$-aspheric.

\printedpage{176}
% head: A. Grothendieck : Dérivateurs, Ch. XII, Édition des Universités de Montpellier II et Paris VII
% foot: Version 30 mars 2003

Finally, suppose that $W$ is \textit{stable under direct factors,} and W(7 bis) satisfied in forward form. I say that then
\[
\mathrm{c}' \ \& \ \mathrm{a}' \Longrightarrow \mathrm{b}',
\]
i.e.
\[
Z_0 \text{ and } Z_3 \ \in \Asph_W \qquad\quad \Longrightarrow \qquad\quad Z_1 \text{ and } Z_2 \ \in \Asph_W.
\]
Factoring $i_1$ and $i_2$ as
\[
\begin{aligned}
Z_0 &\xrightarrow{i'_1} Z'_1 \xrightarrow{\alpha_1} Z_1, \\
Z_0 &\xrightarrow{i'_2} Z'_2 \xrightarrow{\alpha_2} Z_2,
\end{aligned}
\]
with $i'_1$, $i'_2$ open Dwyer immersions, and $\alpha_1, \alpha_2 \in W$, we are reduced to the case where

\msp{page 271}

$i_1$, $i_2$ are \textit{open Dwyer immersions,} and where the square is cocartesian (which implies that it is $W$-cocartesian). Then making the cobase change by $Z_0 \xrightarrow{p_0} Z'_0 = e$, we find a homomorphism of cocartesian squares
\[
\begin{tikzcd}[ampersand replacement=\&, column sep=small, row sep=small]
 \& Z_1 \arrow[rr] \arrow[dd, "p_1" near end] \& \& Z_3 \arrow[dd, "p_3"] \\
Z_0 \arrow[ur] \arrow[rr, crossing over] \arrow[dd, "p_0"'] \& \& Z_2 \arrow[ur] \& \\
 \& Z'_1 \arrow[rr] \& \& Z'_3, \\
e = Z'_0 \arrow[ur] \arrow[rr] \& \& Z'_2 \arrow[ur] \arrow[from=uu, crossing over, "p_2"' near end] \&
\end{tikzcd}
\]
where all the faces of the cube are $W$-cocartesian, and since $p_0 \in W$ ($Z_0$ being $W$-aspheric), the other vertical arrows $p_1$, $p_2$, $p_3$ are \texttt{[in]} $W$, hence the square $Z'_*$ also satisfies $Z'_0$ and $Z'_3$ \texttt{[}$W$\texttt{-]}aspheric (and in addition $Z'_0 = e$), and we must prove $Z'_1$, $Z'_2$ $W$-aspheric (whence likewise for $Z_1$, $Z_2$). Thus, we are reduced to the particular case where $Z_0 = e$.

But in this case there exists a retraction (trivial!) of $Z_1$ onto $Z_0$, hence also of $Z_3$ onto $Z_2$. Thus, $Z_2$ is a direct

\msp{page 272}

factor of a $W$-aspheric object, which implies, by the hypothesis on $W$, that $Z_2$ is $W$-aspheric. By symmetry, $Z_1$ is too. We have won!

\textbf{NB} In fact, these arguments are general, and give this:

\setcounter{footnote}{136}
\textbf{Theorem.} \textit{Let $(\mathcal{M}, W)$ be a model category, admitting a theory of the $W$-homotopy cofiber \textup{(}\footnote{But we only know how to verify the existence of such a theory in $(\mathrm{Cat}, W)$ if we accept the weak form of W(7), i.e. $f_0, f_1, f_2 \in W \implies f_3 \in W$, or some special case of it. But \textit{Brown's lemma} saves the situation, cf. XV \S3.}\textup{)} (i.e. we have a factorization of every arrow $f : X \longrightarrow Y$ as $pi$, with $i \in \Cof_W$, $p \in W$.) We therefore have a notion of $W$-cocartesian square in $\mathcal{M}$ \textup{(}\footnote{to check that the factorization suffices, even if not all objects are fibrant over $e$. But we can then replace $\mathcal{M}$ by $\mathcal{M}_f$, and develop this formalism of homotopy fibers accordingly \dots}\textup{)}. Let
\printedpage{177}
% head: A. Grothendieck : Dérivateurs, Ch. XII, Édition des Universités de Montpellier II et Paris VII
% foot: Version 30 mars 2003
then $f : X \longrightarrow Y$ be a homomorphism from a $W$-cocartesian square $X$ to another $Y$, $f = (f_0, f_1, f_2, f_3)$.}
\[
\begin{tikzcd}[ampersand replacement=\&, column sep=small, row sep=small]
 \& X_2 \arrow[rr] \arrow[dd, "f_2" near end] \& \& X_3 \arrow[dd, "f_3"] \\
X_0 \arrow[ur] \arrow[rr, crossing over] \arrow[dd, "f_0"'] \& \& X_1 \arrow[ur] \& \\
 \& Y_2 \arrow[rr] \& \& Y_3 \\
Y_0 \arrow[ur] \arrow[rr] \& \& Y_1 \arrow[ur] \arrow[from=uu, crossing over, "f_1"' near end] \&
\end{tikzcd}
\]

\textit{Let us consider the conditions on $f$}

\begin{enumerate}
\item[\textit{a)}] \textit{$f_0 \in W$.}
\item[\textit{b)}] \textit{$f_1$ and $f_2 \in W$.}
\item[\textit{c)}] \textit{$f_3 \in W$.}
\end{enumerate}

\textit{Then}

\begin{enumerate}
\item[\textit{1°)}] \textit{a \& b $\Longrightarrow$ c.}

\msp{page 273}

\item[\textit{2°)}] \textit{Suppose that $\mathcal{M}$ has a final object $e$. Then the implication}
\[
\textit{b \& c} \Longrightarrow \textit{a}
\]
\textit{for}~every \textit{homomorphism $f : X \longrightarrow Y$ of $W$-cocartesian squares is equivalent to the following axiom on $W$:}

Axiom of conservative $W$-cocartesian squares (\footnote{cf. argument page 78 which shows that the implication b \& c $\Longrightarrow$ a entails the conservative axiom.}): \textit{If}
\[
\begin{tikzcd}[ampersand replacement=\&, column sep=small, row sep=small]
 \& Z_1 \arrow[dr, "j_2"] \& \\
Z_0 \arrow[ur, "i_1"] \arrow[dr, "i_2"'] \& \& Z_3 \\
 \& Z_2 \arrow[ur, "j_1"'] \&
\end{tikzcd}
\]
\textit{is $W$-cocartesian, then $j_1 \in W \Longrightarrow i_1 \in W$ (hence in fact $i_1 \in W \Longleftrightarrow j_1 \in W$ and $i_2 \in W \Longleftrightarrow j_2 \in W$).}

\textit{This axiom is equivalent to the conjunction of the following two axioms.}

\textbf{Cons 1.} \textit{If we have a cocartesian square}
\printedpage{178}
% head: A. Grothendieck : Dérivateurs, Ch. XII, Édition des Universités de Montpellier II et Paris VII
% foot: Version 30 mars 2003
\[
\begin{tikzcd}[ampersand replacement=\&, column sep=small, row sep=small]
S \arrow[rr] \arrow[dr] \& \& T \arrow[dr] \& \\
 \& e \arrow[rr] \& \& U,
\end{tikzcd}
\]
\textit{with $f \in \Cof_W$, then $U \in \Asph_W \Longrightarrow f \in W$ (hence $U \in \Asph_W \Longleftrightarrow f \in W$).}

\textbf{Cons 2.} \textit{If we have a cartesian square,}
\[
\begin{tikzcd}[ampersand replacement=\&, column sep=small, row sep=small]
 \& Z_1 \arrow[dr] \& \\
Z_0 \arrow[ur, "i_1"] \arrow[dr, "i_2"'] \& \& Z_3, \\
 \& Z_2 \arrow[ur] \&
\end{tikzcd}
\]
\textit{with $i_1, i_2 \in \Cof_W$, then: If $Z_1$, $Z_2$, $Z_3$ are $W$-aspheric, $Z_0$ is too.}

\msp{page 274}

\item[\textit{3°)}] \textit{Suppose that every direct factor of a $W$-aspheric object is $W$-aspheric, which is the case in particular if $W$ itself is stable under direct factors, a fortiori if $W$ is strongly saturated. Suppose in addition that $\mathcal{M}$ has} \texttt{[a]} \textit{final object. Then we have (under the general conditions at the beginning)}
\[
\textit{c \& a} \Longrightarrow \textit{b,}
\]
\textit{i.e. if $f_0$, $f_3$ are} \texttt{[in]} \textit{$W$, likewise $f_1$, $f_2$.}
\end{enumerate}

This last part can be unfolded as follows:

\textbf{Corollary 1.} \textit{Under the additional hypothesis} \textup{3\,°)} \textit{on $W$, the validity of c \& a $\Longrightarrow$ b (for every arrow $f$ between $W$-cocartesian squares) is equivalent to each of the following two other axioms on $W$ (which are therefore equivalent):}

\begin{enumerate}
\item[\textit{1)}] \textit{If we have a $W$-cocartesian square}
\[
\begin{tikzcd}[ampersand replacement=\&, column sep=small, row sep=small]
 \& Z_1 \arrow[dr] \& \\
Z_0 \arrow[ur] \arrow[dr] \& \& Z_3, \\
 \& Z_2 \arrow[ur] \&
\end{tikzcd}
\]
\textit{then $Z_0$ and $Z_3$ $W$-aspheric $\Longrightarrow$ $Z_1$ and $Z_2$ $W$-aspheric.}
\item[\textit{1$'$)}] \textit{As in 1), with $Z_0 = e$ (in the case where we suppose that $\mathcal{M}$ admits} \texttt{[a]} \textit{final object $e$).}
\end{enumerate}

\msp{page 275}

\textbf{Corollary 2.} \textit{Suppose that $W$ satisfies the axiom of conservative $W$-cocartesian squares, that every direct factor of a $W$-aspheric object is $W$-aspheric (e.g. $W$ stable under direct factors), finally that $\mathcal{M}$ admits a final object. Then for every homomorphism $f : X \longrightarrow Y$ of $W$-cocartesian squares, the conjunction of two of the conditions a) b) c) on p. 272 implies the third.} % [sic] pout -> pour, typographical error

\end{document}
